\begin{enumerate}
\item Biểu diễn tập hợp $A$ trên trục số (gạch bỏ phần ngoài $A$).
\item Biểu diễn tập hợp $B$ và gạch bỏ toàn bộ phần thuộc $B$; phần còn lại không bị gạch là $A \setminus B$.
\item Với phần bù $C_{\mathbb{R}} A = \mathbb{R} \setminus A$: bỏ đi phần thuộc $A$ trên trục số thực $\mathbb{R}$.
\end{enumerate}
Lưu ý: Khi gạch bỏ một khoảng mở $(a; b)$, các đầu mút $a, b$ vẫn được giữ lại (đổi thành ngoặc vuông); ngược lại, khi gạch bỏ đoạn $[a; b]$, các đầu mút $a, b$ bị loại đi (đổi thành ngoặc tròn).
